\newpage
\section{Resultados esperados}
\label{sc:resultados}

Ao término do TCC02, espera-se que o Dança~AI comprove a viabilidade técnica do pipeline proposto em hardware de
consumo e a eficácia do feedback automatizado sobre indicadores mensuráveis de execução nos
quatro módulos. Os resultados esperados estão organizados em quatro
dimensões, alinhadas aos objetivos específicos 3 a 7.

\subsection{Protótipo Funcional do Sistema Dança~AI}
\label{sub:res_prototipo}

Espera-se a entrega de um aplicativo Android funcional que implemente integralmente os
quatro módulos de análise descritos nos objetivos 3 a 6: captura de vídeo via CameraX e
extração em tempo real dos 33~\textit{landmarks} corporais pelo MediaPipe Pose Landmarker;
módulo de transferência de peso, classificando o apoio como esquerdo, direito ou neutro com
\textit{feedback} visual imediato; módulo de postura, avaliando o alinhamento do tronco e dos
ombros por regras geométricas e classificador TFLite embarcado; módulo de ritmo, com
metrônomo configurável por BPM, \textit{feedback} auditivo e validação rítmica em tempo real;
e módulo de identificação de movimentos, com classificador TFLite para os seis passos do
forró universitário nos modos Livre e Desafio. O aplicativo deverá operar de forma estável
durante sessões contínuas de ao menos 5~minutos sem degradação de desempenho em dispositivos
de segmento intermediário.

A interface, construída em Jetpack Compose, contemplará telas de configuração de sessão ---
com guia de enquadramento, seleção de modo (Livre ou Desafio) e configuração de BPM ---,
treino em tempo real com \textit{overlay} de \textit{feedback} e relatório de sessão com
indicadores por módulo. Esse fluxo permite ao praticante iniciar e encerrar sessões de forma
autônoma, sem supervisão de um instrutor~\cite{stanford_ai_dance_coaching}.

\subsection{Validação do Desempenho Computacional}
\label{sub:res_desempenho}

Espera-se que o \textit{pipeline} completo --- captura, inferência MediaPipe, cálculo de
ângulos, comparação DTW e renderização --- opere com latência mediana inferior a 100~ms em
dispositivos Android de segmento intermediário, atendendo ao orçamento estabelecido na
Seção~\ref{subsub:otimizacao} com base nos limites neurológicos de percepção de
responsividade documentados na literatura~\cite{plos_feedback_delay}. Os dados de latência
serão apresentados por estágio do \textit{pipeline}, com distinção entre o tempo de
inferência do modelo (estimado entre 20--40~ms com \textit{delegate} GPU, conforme Lee et
al.~\cite{lee2019ondevice}) e os demais estágios de processamento. Espera-se ainda taxa
mínima de 20~FPS, suficiente para que o \textit{overlay} seja percebido como fluido durante a
execução dos movimentos.

A acurácia dos ângulos articulares calculados a partir dos \textit{landmarks} do MediaPipe
deverá situar-se dentro dos limites de concordância estabelecidos por Latreche et
al.~\cite{LATRECHE2023112826}, que reportaram valores clinicamente aceitáveis para aplicações
de análise de movimento --- domínio com requisitos de precisão comparáveis aos do Dança~AI,
conforme discutido na Seção~\ref{sub:visao}.

\subsection{Módulos de Classificação e Comparação de Movimentos}
\label{sub:res_dtw}

Espera-se que os classificadores TFLite de postura e de identificação de movimentos
apresentem acurácia global superior a 85\% no conjunto de teste, com precisão e
\textit{recall} por classe documentados para cada um dos seis passos do forró universitário
(Balanço, Base~1, Base~2, Abertura, Giro Invertido e Giro em
Linha)~\cite{roggio2024pose, bursa2023tflite}. O módulo de ritmo deverá detectar desvios
entre a cadência de transferências de peso e o BPM configurado com margem de erro inferior a
um compasso, gerando \textit{feedback} visual coerente com a percepção do
praticante~\cite{rose2019metronome}.

O módulo de comparação via DTW deverá produzir pontuações que discriminem de forma
consistente entre execuções corretas e incorretas, com correlação verificável com a avaliação
qualitativa de um instrutor qualificado de forró universitário~\cite{dynamic_dance_warping}. A
implementação em Kotlin puro com restrição de janela de Sakoe-Chiba deverá apresentar tempo
de execução inferior a 5~ms por sequência, de modo que o módulo não represente gargalo no
orçamento de latência global~\cite{skeleton_dtw_sport}.

\subsection{Eficácia Pedagógica e Usabilidade}
\label{sub:res_pedagogia}

Espera-se que os testes com usuários evidenciem melhoria mensurável na acurácia de execução
após sessões de prática assistida pelo Dança~AI, avaliada pelos índices de conformidade de
cada módulo --- transferência de peso, postura, ritmo e identificação de movimentos --- antes
e após cada sessão. Como referência, o estudo de
Zhou~\cite{stanford_ai_dance_coaching} documentou ganho de 55,9\% para 66,8\% de acurácia
em grupo de dez praticantes com \textit{feedback} baseado em estimativa de pose, com 85\% das
intervenções do sistema avaliadas como coerentes pelos participantes. Espera-se que os dados
obtidos no presente experimento confirmem tendência de melhoria análoga no contexto específico
do forró universitário, eventualmente com variações associadas ao nível de experiência prévia
dos participantes.

Para a dimensão de usabilidade, espera-se pontuação SUS média superior a 70 --- limiar
correspondente à faixa de usabilidade classificada como ``boa'' --- indicando que o sistema é
percebido como utilizável por praticantes sem formação técnica. As entrevistas semiestruturadas
deverão fornecer evidências qualitativas sobre quais dimensões do \textit{feedback} ---
transferência de peso, postural, rítmica ou de identificação de movimentos --- são
percebidas como mais úteis pelos praticantes, orientando a priorização de melhorias para
versões futuras do sistema. Em conjunto, esses resultados permitirão posicionar o Dança~AI
como contribuição original ao estado da arte sintetizado na
Tabela~\ref{tab:correlatos}: um sistema que reúne, pela primeira vez em um único aplicativo
autônomo, inferência completamente \textit{on-device}, avaliação biomecânica baseada em
ângulos articulares, comparação de sequências via DTW e foco em uma modalidade de dança de
salão específica~\cite{THARATIPYAKUL2024e36589}.
